\BourbakiExerciseGroup

\begin{enumerate}
\def\labelenumi{\arabic{enumi})}
\tightlist
\item
  多项式 \(X^{2}-2\) 在 Q{[}X{]} 中不可约（参见 I，第 51 页，定理 7）；令 \(\alpha\) 是它在 Q 的一个代数闭包 \(\Omega\) 中的一个根。证明多项式 \(X^{2}-\alpha\) 在 \(E = Q(\alpha)\) 上不可约；设 \(\beta\) 是该多项式在 \(\Omega\) 中的一个根，并且令 \(F = E(\beta) = Q(\beta)\) 。证明 F 不是 Q 的拟伽罗瓦扩张；由 F 生成的 Q 的拟伽罗瓦扩张是什么？
\end{enumerate}

* 2) 设 E 为 Q 在 R 中的相对代数闭包（“实代数数域”）；若 u 是 E 的一个自同构，则 \(u(x) = x\) 对所有 \(x \in \mathbf{Q}\) 成立；证明若 \(y \in E\) 使得 \(y > 0\) ，则也有 \(u(y) > 0\) 。由此推出 \(u(y) = y\) 对所有 \(y \in E\) 成立. *

\begin{enumerate}
\def\labelenumi{\arabic{enumi})}
\setcounter{enumi}{2}
\item
  证明域 \(\mathbf{K}\) 的每个代数扩张，若它由一组在 \(\mathbf{K}\) 上次数为 2 的元素生成，则在 \(\mathbf{K}\) 上是拟伽罗瓦的.
\item
  设 K 为一个域，f 为 K{[}X{]} 中在 K 上可分且次数为 n 的不可约多项式，并设 \(\alpha_{i}\) ( \(1 \leqslant i \leqslant n\) ) 为其在 K 的一个代数闭包 \(\Omega\) 中的根。设 g 为 K{[}X{]} 中的一个多项式，且 h 为 K{[}X{]} 中多项式 \(f(g(\mathbf{X}))\) 的一个不可约因子。证明 h 的次数是 n 的倍数 rn，并且 h 与每个多项式 \(g(\mathbf{X}) - \alpha_{i}\) ( \(1 \leqslant i \leqslant n\) )（属于 \(\Omega[\mathbf{X}]\) ）恰好有 r 个公共根（考虑 h 的一个根在 K 上的共轭）。
\item
  \begin{enumerate}
  \def\labelenumii{\alph{enumii})}
  \tightlist
  \item
    设 K 为一个域， \(\Omega\) 为 K 的一个代数闭包，且 E、F 为 \(\Omega\) 的两个子扩张。证明 E 与 F 的每个复合扩张（V，第 12 页）都同构于形如 \((L_{w}, 1, w)\) 的一个复合扩张，其中 w 表示 F 到 \(\Omega\) 的一个子域上的 K-同构，1 是 E 的恒等映射，且 \(L_{w}\) 是 \(\Omega\) 中由 \(E \cup w(F)\) 生成的子域.
  \end{enumerate}
\end{enumerate}

\begin{enumerate}
\def\labelenumi{\alph{enumi})}
\setcounter{enumi}{1}
\item
  由 a) 推出，若 F 在 K 上的次数为有限的 n，则 E 与 F 的复合扩张中两两不同构的至多有 n 个。
\item
  设 E 是 K 的拟伽罗瓦扩张，F 是 E 的一个子扩张。证明 E 与 F 的每个复合扩张都同构于形如 \((E, 1, v)\) 的一个复合扩张，其中 v 是 F 到 E 的一个子域上的 K-同构。
\end{enumerate}

\begin{enumerate}
\def\labelenumi{\arabic{enumi})}
\setcounter{enumi}{5}
\tightlist
\item
  \begin{enumerate}
  \def\labelenumii{\alph{enumii})}
  \tightlist
  \item
    设 K 是特征指数为 p 的域，L 是 K 的一个扩张， \(\Omega\) 为 L 的一个代数闭包，以及 \(\bar{K}\) 是 K 在 \(\Omega\) 中的相对代数闭包。证明 \(\bar{K}\) 与 \(L^{p^{-\infty}}\) 在 \(E = \bar{K} \cap L^{p^{-\infty}}\) 上线性不相交。（设 \((a_{\lambda})\) 是向量空间 \(L^{p^{-\infty}}\) 在 E 上的一个基；考虑一个关系 \(\sum_{\lambda} c_{\lambda} a_{\lambda} = 0\) ，其中 \(c_{\lambda} \in \bar{K}\) 且非零的
  \end{enumerate}
\end{enumerate}

\(c_{\lambda}\) 的个数尽可能小，并把 \(\Omega\) 的一个 L-自同构作用于这个关系上。）

\begin{enumerate}
\def\labelenumi{\alph{enumi})}
\setcounter{enumi}{1}
\tightlist
\item
  由a)推出，若E和F是K的两个扩张，使得E在K上是代数的，且K在F中是相对代数闭的，则E和F的任意两个复合扩张都是同构的。
\end{enumerate}

\begin{enumerate}
\def\labelenumi{\arabic{enumi})}
\setcounter{enumi}{6}
\tightlist
\item
  设 K 是特征为 0 的域， \(\alpha\) , \(\beta\) , \(\xi\) 是 K 的一个代数闭包中的三个元素，使得 \([\mathbf{K}(\alpha):\mathbf{K}] = 2\) , \([\mathbf{K}(\beta):\mathbf{K}] = 3\) , \([\mathbf{K}(\xi):\mathbf{K}] = n > 1\) ；用 \(\alpha'\) （相应地 \(\beta'\) , \(\beta''\) ) 表示唯一的共轭元 \(\neq \alpha\) （ \(\alpha\) 的）（相应地，那些共轭元 \(\neq \beta\) （ \(\beta\) 的））.
\end{enumerate}

\begin{enumerate}
\def\labelenumi{\alph{enumi})}
\item
  证明： \([\mathbf{K}(\alpha, \xi): \mathbf{K}] = n\) 如果 \(\alpha \in \mathbf{K}(\xi)\) ，否则 \([\mathbf{K}(\alpha, \xi): \mathbf{K}] = 2n\) 。
\item
  证明： \([\mathbf{K}(\beta,\xi):\mathbf{K}]=n\) 如果 \(\beta\in\mathbf{K}(\xi)\) ；\([\mathbf{K}(\beta,\xi):\mathbf{K}]=2n\) 如果 \(\beta\in\mathbf{K}(\xi)\) 并且 \(\beta'\) , \(\beta''\) 中恰有一个属于 \(\mathbf{K}(\xi)\) ；最后 \([\mathbf{K}(\beta,\xi):\mathbf{K}]=3n\) 如果三个元素 \(\beta\) , \(\beta'\) , \(\beta''\) 中没有一个属于 \(\mathbf{K}(\xi)\) .
\item
  如果 \(d\) 是 \(\beta\) 的极小多项式的判别式，证明 \(\mathbf{K}(\beta, \beta') = \mathbf{K}(\beta, \sqrt{d})\) .
\end{enumerate}

\(d\) ) 证明 \(\mathbf{K}(\beta, \beta') = \mathbf{K}(\beta - \beta')\) .

\begin{enumerate}
\def\labelenumi{\alph{enumi})}
\setcounter{enumi}{4}
\item
  如果 \(n = 2\) 且 \(\alpha \notin \mathbf{K}(\xi)\) ，证明 \(\mathbf{K}(\alpha, \xi) = \mathbf{K}(\alpha + \xi)\) .
\item
  证明： \(\mathbf{K}(\alpha, \beta) = \mathbf{K}(\alpha + \beta)\) .
\end{enumerate}

\(g\) ) 证明 \(\mathbf{K}(\alpha ,\beta) = \mathbf{K}(\alpha \beta)\) ，除非 \(\alpha = aj\) 且 \(\beta^3 = b\) ，其中 \(a,b\) 是 \(\mathbf{K}\) 的元素且 \(j^3 = 1\) .

\begin{enumerate}
\def\labelenumi{\arabic{enumi})}
\setcounter{enumi}{7}
\tightlist
\item
  设 K 为无限域，E 为 K 的有限次代数扩张，且 \(\tau\) 是向量 K-空间 E 的一个自同构，使得对所有 \(x \in E\) , \(\tau(x)\) 在 K 上与 x 共轭。证明 \(\tau\) 是域 E 的一个 K-自同构。（设 N 为由 E 生成的 K 的拟伽罗瓦扩张，并设 \(\Gamma\) 为 N 的 K-自同构群， \(\Delta\) 为 \(\Gamma\) 的由所有 E-自同构组成的子群，并且 \(\sigma_{j}\) ( \(1 \leqslant j \leqslant r\) ) 是 \(\Delta\) 的左陪集的一个代表系。若 \((u_{i})_{1 \leqslant i \leqslant n}\) 是向量 K-空间 E 的一组基，考虑多项式
\end{enumerate}

\[
\mathrm{F} \left(\mathrm{X} _ {1}, \dots , \mathrm{X} _ {n}\right) = \prod_ {j = 1} ^ {r} \left(\sum_ {i = 1} ^ {n} \left(\tau \left(u _ {i}\right) - \sigma_ {j} \left(u _ {i}\right)\right) \mathrm{X} _ {i}\right).
\]

